% Options for packages loaded elsewhere
% Options for packages loaded elsewhere
\PassOptionsToPackage{unicode}{hyperref}
\PassOptionsToPackage{hyphens}{url}
\PassOptionsToPackage{dvipsnames,svgnames,x11names}{xcolor}
%
\documentclass[
  french,
  10pt,
  a4paper,
]{scrartcl}
\usepackage{xcolor}
\usepackage[margin=2cm]{geometry}
\usepackage{amsmath,amssymb}
\setcounter{secnumdepth}{5}
\usepackage{iftex}
\ifPDFTeX
  \usepackage[T1]{fontenc}
  \usepackage[utf8]{inputenc}
  \usepackage{textcomp} % provide euro and other symbols
\else % if luatex or xetex
  \usepackage{unicode-math} % this also loads fontspec
  \defaultfontfeatures{Scale=MatchLowercase}
  \defaultfontfeatures[\rmfamily]{Ligatures=TeX,Scale=1}
\fi
\usepackage{lmodern}
\ifPDFTeX\else
  % xetex/luatex font selection
\fi
% Use upquote if available, for straight quotes in verbatim environments
\IfFileExists{upquote.sty}{\usepackage{upquote}}{}
\IfFileExists{microtype.sty}{% use microtype if available
  \usepackage[]{microtype}
  \UseMicrotypeSet[protrusion]{basicmath} % disable protrusion for tt fonts
}{}
\makeatletter
\@ifundefined{KOMAClassName}{% if non-KOMA class
  \IfFileExists{parskip.sty}{%
    \usepackage{parskip}
  }{% else
    \setlength{\parindent}{0pt}
    \setlength{\parskip}{6pt plus 2pt minus 1pt}}
}{% if KOMA class
  \KOMAoptions{parskip=half}}
\makeatother
% Make \paragraph and \subparagraph free-standing
\makeatletter
\ifx\paragraph\undefined\else
  \let\oldparagraph\paragraph
  \renewcommand{\paragraph}{
    \@ifstar
      \xxxParagraphStar
      \xxxParagraphNoStar
  }
  \newcommand{\xxxParagraphStar}[1]{\oldparagraph*{#1}\mbox{}}
  \newcommand{\xxxParagraphNoStar}[1]{\oldparagraph{#1}\mbox{}}
\fi
\ifx\subparagraph\undefined\else
  \let\oldsubparagraph\subparagraph
  \renewcommand{\subparagraph}{
    \@ifstar
      \xxxSubParagraphStar
      \xxxSubParagraphNoStar
  }
  \newcommand{\xxxSubParagraphStar}[1]{\oldsubparagraph*{#1}\mbox{}}
  \newcommand{\xxxSubParagraphNoStar}[1]{\oldsubparagraph{#1}\mbox{}}
\fi
\makeatother


\usepackage{longtable,booktabs,array}
\usepackage{calc} % for calculating minipage widths
% Correct order of tables after \paragraph or \subparagraph
\usepackage{etoolbox}
\makeatletter
\patchcmd\longtable{\par}{\if@noskipsec\mbox{}\fi\par}{}{}
\makeatother
% Allow footnotes in longtable head/foot
\IfFileExists{footnotehyper.sty}{\usepackage{footnotehyper}}{\usepackage{footnote}}
\makesavenoteenv{longtable}
\usepackage{graphicx}
\makeatletter
\newsavebox\pandoc@box
\newcommand*\pandocbounded[1]{% scales image to fit in text height/width
  \sbox\pandoc@box{#1}%
  \Gscale@div\@tempa{\textheight}{\dimexpr\ht\pandoc@box+\dp\pandoc@box\relax}%
  \Gscale@div\@tempb{\linewidth}{\wd\pandoc@box}%
  \ifdim\@tempb\p@<\@tempa\p@\let\@tempa\@tempb\fi% select the smaller of both
  \ifdim\@tempa\p@<\p@\scalebox{\@tempa}{\usebox\pandoc@box}%
  \else\usebox{\pandoc@box}%
  \fi%
}
% Set default figure placement to htbp
\def\fps@figure{htbp}
\makeatother



\ifLuaTeX
\usepackage[bidi=basic]{babel}
\else
\usepackage[bidi=default]{babel}
\fi
% get rid of language-specific shorthands (see #6817):
\let\LanguageShortHands\languageshorthands
\def\languageshorthands#1{}


\setlength{\emergencystretch}{3em} % prevent overfull lines

\providecommand{\tightlist}{%
  \setlength{\itemsep}{0pt}\setlength{\parskip}{0pt}}



 
\usepackage[numbers,round,sort,compress]{natbib}
\bibliographystyle{unsrturl}


% Configuration PDF commune aux documents Husson.
\PassOptionsToPackage{french}{babel}
\PassOptionsToPackage{bidi=default}{babel}
\usepackage{fontspec}
\usepackage{listings}
\defaultfontfeatures{SmallCapsFeatures={Letters=SmallCaps}}
\IfFontExistsTF{texgyretermes-regular.otf}{%
  \defaultfontfeatures[Times New Roman]{%
    SmallCapsFont = texgyretermes-regular.otf,%
    SmallCapsFeatures = {Letters=SmallCaps}%
  }%
}{}
\usepackage{xcolor}
\usepackage{fvextra}
\usepackage{newunicodechar}
\usepackage{enumitem}
\usepackage{etoolbox}
\usepackage{needspace}
\usepackage{pdflscape}
\usepackage{adjustbox}
\usepackage{ragged2e}

% Securite et compatibilite pour documents existants (classes KOMA vs standard)
\providecommand{\setkomafont}[2]{}
\providecommand{\addtokomafont}[2]{}
\providecommand{\RedeclareSectionCommand}[2][]{}

% Transition automatique de numerotation : chiffres romains pour les liminaires (TOC), puis reprise en chiffres arabes des le corps du texte
\AddToHook{cmd/tableofcontents/after}{%
  \clearpage
  \pagenumbering{arabic}%
  \setcounter{page}{1}%
}

\definecolor{headingblue}{HTML}{17365D}
\definecolor{headingred}{HTML}{7A1F2B}
\definecolor{headinggreen}{HTML}{1F5A3A}
\definecolor{headingcharcoal}{HTML}{34383D}
\definecolor{commentgray}{HTML}{6F6F6F}
\definecolor{pythonaccent}{HTML}{9A6100}
\definecolor{pythonfill}{HTML}{FFF5D9}
\definecolor{raccent}{HTML}{276DC3}
\definecolor{rfill}{HTML}{EAF2FB}

\newunicodechar{≥}{\ensuremath{\geq}}
\newunicodechar{≤}{\ensuremath{\leq}}
\newunicodechar{≈}{\ensuremath{\approx}}
\newunicodechar{₂}{\textsubscript{2}}
\newunicodechar{₁}{\textsubscript{1}}
\newunicodechar{₀}{\textsubscript{0}}
\newunicodechar{₃}{\textsubscript{3}}
\newunicodechar{₄}{\textsubscript{4}}
\newunicodechar{ρ}{\ensuremath{\rho}}
\setlength{\emergencystretch}{3em}

% Typographie française : "Tableau" au lieu de "Table" lorsque la langue principale est le français
\makeatletter
\AddToHook{begindocument/end}{%
  \@ifundefined{languagename}{}{%
    \def\@frenchname{french}%
    \ifx\languagename\@frenchname
      \def\tablename{Tableau}%
    \fi
  }%
}
\makeatother


% Pagination éditoriale. Une ligne isolée au début ou à la fin d'une page est
% évitée, tandis que \raggedbottom autorise un blanc final raisonnable plutôt
% que d'étirer artificiellement les espacements verticaux.
\clubpenalty=10000
\widowpenalty=10000
\displaywidowpenalty=10000
\raggedbottom

\makeatletter
\renewcommand*{\@pnumwidth}{2.8em}
\renewcommand*{\@tocrmarg}{3.6em}
\makeatother

% Hiérarchie éditoriale sobre : les trois premiers niveaux structurent la page
% par la graisse et une couleur foncée. Les niveaux suivants restent proches
% du texte courant et se distinguent par la graisse, l'italique ou les petites
% capitales. Seuls les trois premiers niveaux sont numérotés.
\setcounter{secnumdepth}{3}
\setkomafont{disposition}{\rmfamily}
\setkomafont{section}{\rmfamily\color{headingblue}\bfseries\Large}
\setkomafont{subsection}{\rmfamily\color{headingred}\bfseries\large}
\setkomafont{subsubsection}{\rmfamily\color{headinggreen}\bfseries\normalsize}
\setkomafont{paragraph}{\rmfamily\color{headingcharcoal}\bfseries\normalsize}
\setkomafont{subparagraph}{\rmfamily\color{headingcharcoal}\itshape\normalsize}

\RedeclareSectionCommand[
  beforeskip=1.25\baselineskip,
  afterskip=0.65\baselineskip
]{section}
\RedeclareSectionCommand[
  beforeskip=1.05\baselineskip,
  afterskip=0.50\baselineskip
]{subsection}
\RedeclareSectionCommand[
  beforeskip=0.90\baselineskip,
  afterskip=0.42\baselineskip
]{subsubsection}
\RedeclareSectionCommand[
  beforeskip=0.76\baselineskip,
  afterskip=0.28\baselineskip
]{paragraph}
\RedeclareSectionCommand[
  beforeskip=0.62\baselineskip,
  afterskip=0.24\baselineskip
]{subparagraph}

% Pandoc rend normalement un titre Markdown de niveau 6 comme du texte brut
% en LaTeX. Le filtre heading-levels.lua l'entoure de cet environnement pour
% lui donner un style discret et un espace propre avec le texte.
\newenvironment{hussonsixthheading}{%
  \Needspace{3\baselineskip}%
  \par\addvspace{0.52\baselineskip}%
  \noindent\rmfamily\color{headingcharcoal}\scshape\small
}{%
  \par\nobreak\vspace{0.20\baselineskip}%
}

\setlength{\parindent}{0pt}
\setlength{\parskip}{3pt}
\setlist[itemize]{itemsep=1pt,parsep=1pt,topsep=2pt}
\setlist[enumerate]{itemsep=1pt,parsep=1pt,topsep=2pt}

% Espacement vertical compact et sobre pour les tableaux (hauteur minimale).
% La boîte typographique normale est rétablie avant de composer les paragraphes
% multiligne : la marge n'est donc appliquée qu'une fois par cellule.
\renewcommand{\arraystretch}{1}
\setlength{\extrarowheight}{0pt}
\makeatletter
\newbox\hussontablestrutbox
\newcommand{\hussontablepadding}{%
  \setbox\hussontablestrutbox=\copy\strutbox
  \setbox\strutbox=\hbox{%
    \vrule width 0pt
      height \dimexpr\ht\hussontablestrutbox+2.5pt\relax
      depth \dimexpr\dp\hussontablestrutbox+2.5pt\relax
  }%
  \let\hussontablesavedmkpream\@mkpream
  \def\@mkpream##1{%
    \hussontablesavedmkpream{##1}%
    \setbox\strutbox=\copy\hussontablestrutbox
    \let\@mkpream\hussontablesavedmkpream
  }%
}
\makeatother

% Les tableaux gardent leur hauteur naturelle, avec une taille homogène. Un
% tableau long peut se poursuivre sur la page suivante, mais ne commence pas
% dans les toutes dernières lignes de la page courante.
\AtBeginEnvironment{longtable}{\Needspace{6\baselineskip}\small\hussontablepadding}
\AtBeginEnvironment{tabular}{\small\hussontablepadding}

% Filet de securite : seuls les tableaux depassant la ligne sont reduits a \linewidth.
\BeforeBeginEnvironment{tabular*}{\begin{adjustbox}{max width=\linewidth}}
\AfterEndEnvironment{tabular*}{\end{adjustbox}}
\BeforeBeginEnvironment{tabular}{\begin{adjustbox}{max width=\linewidth}}
\AfterEndEnvironment{tabular}{\end{adjustbox}}

% Autorise la coupure et la césure automatique des mots dans les cellules de tableau :
% Le \raggedright standard de LaTeX désactive toute césure (\hyphenpenalty=10000).
% \RaggedRight rétablit les règles de césure automatique de la langue (français),
% et \hspace{0pt} permet au moteur TeX de couper automatiquement dès le premier mot d'une cellule.
\AtBeginDocument{
  \renewcommand{\raggedright}{\RaggedRight\hspace{0pt}}
}

% Les blocs de code restent divisibles lorsqu'ils sont longs. Cette réserve
% empêche seulement qu'un bloc commence en bas de page avec une ou deux lignes.
\BeforeBeginEnvironment{Shaded}{\Needspace{8\baselineskip}}

% Les blocs de code sont compacts et peuvent revenir à la ligne.
\DefineVerbatimEnvironment{Highlighting}{Verbatim}{%
  fontsize=\footnotesize,
  breaklines,
  breakanywhere,
  commandchars=\\\{\}
}
\providecommand{\CommentTok}[1]{\textcolor{commentgray}{\textit{#1}}}

\newcommand{\codelanguagelabel}[2]{%
  \Needspace{5\baselineskip}%
  \par\addvspace{0.55\baselineskip}%
  \noindent\begingroup
  \setlength{\fboxsep}{3pt}%
  \colorbox{#1}{\color{white}\sffamily\bfseries\scriptsize #2}%
  \endgroup\par\nobreak\vspace{-0.25\baselineskip}%
}

\newenvironment{pythoncode}{%
  \codelanguagelabel{pythonaccent}{PYTHON}%
  \begingroup\colorlet{shadecolor}{pythonfill}%
}{%
  \endgroup\par\addvspace{0.25\baselineskip}%
}

\newenvironment{rcode}{%
  \codelanguagelabel{raccent}{R}%
  \begingroup\colorlet{shadecolor}{rfill}%
}{%
  \endgroup\par\addvspace{0.25\baselineskip}%
}

\makeatletter
\renewcommand{\maketitle}{
  \begin{center}
    {\Large\bfseries\rmfamily \@title\par}
    \vskip 0.35em
    {\large\rmfamily \@subtitle\par}
    \vskip 0.65em
    {\normalsize\rmfamily \@author\par}
  \end{center}
}
\makeatother

\renewcommand{\contentsname}{}
\AtBeginDocument{
  \let\oldtableofcontents\tableofcontents
  \renewcommand{\tableofcontents}{
    \begingroup
    \footnotesize
    \setlength{\parskip}{1pt}
    \oldtableofcontents
    \endgroup
  }
}
% navigation-header.tex
% Module LaTeX pour la barre de navigation hierarchique sur une seule ligne
% Extension Husson pour Quarto PDF

\RequirePackage{scrlayer-scrpage}
\RequirePackage{xcolor}
\RequirePackage{graphicx}
\RequirePackage{geometry}

% Calibrage de l'en-tete et du pied de page pour eviter tout chevauchement et avertissement KOMA-Script
\geometry{headheight=18pt, headsep=14pt, footskip=28pt}

% Couleurs complementaires
\providecolor{headingblue}{HTML}{17365D}
\providecolor{headingcharcoal}{HTML}{34383D}
\providecolor{commentgray}{HTML}{6F6F6F}
\providecolor{linegray}{HTML}{D0D4D9}

% Elements graphiques de separation et d'accentuation
\makeatletter
\newdimen\husson@w
\newsavebox{\husson@fitbox}

% Branche active
\newcommand{\hussonactive}[1]{#1}

% Separateur de fil d'Ariane
\newcommand{\hussonsep}{\,\textcolor{commentgray!60}{›}\,}

% Espacement elastique entre les rubriques H2
\newcommand{\hussonhsep}{\hspace{1.1em plus 0.8em minus 0.5em}}

% Separateur vertical entre H1 et la liste des H2
\newcommand{\hussonbarsep}{\hspace{0.8em}\textcolor{linegray}{\rule[-0.2em]{0.6pt}{1.0em}}\hspace{0.8em}}

% H1 autonome (chapitre sans H2 : Bibliographie, Conclusion, Annexes, etc.)
\newcommand{\hussonnavstandalone}[1]{%
  \noindent\normalfont\scriptsize
  \textcolor{headingblue}{\textbf{#1}}%
}

% Ajustement automatique doux : si la ligne depasse \textwidth, reduction proportionnelle
\newcommand{\hussonnavfit}[1]{%
  \sbox{\husson@fitbox}{#1}%
  \ifdim\wd\husson@fitbox>\textwidth
    \resizebox{\textwidth}{!}{\usebox{\husson@fitbox}}%
  \else
    \usebox{\husson@fitbox}%
  \fi
}
\makeatother

% Gestion robuste des marques TeX et du pied de page (LaTeX 2022+ / Expl3)
\ExplSyntaxOn
\NewMarkClass{hussonnav}

\tl_new:N \g_husson_title_tl
\tl_new:N \g_husson_author_tl
\tl_new:N \g_husson_footerleft_tl

\NewDocumentCommand{\hussondocauthor}{m}{
  \tl_gset:Nn \g_husson_author_tl {#1}
}
\NewDocumentCommand{\hussondoctitle}{m}{
  \tl_gset:Nn \g_husson_title_tl {#1}
}
\NewDocumentCommand{\hussondocfooterleft}{m}{
  \tl_gset:Nn \g_husson_footerleft_tl {#1}
}

\cs_new:Npn \husson_set_state:nnn #1 #2 #3 {
  \tl_gset:cn { g_husson_h_#1_tl } { #2 }
  \tl_gset:cn { g_husson_bar_#1_tl } { #3 }
}

\cs_new:Npn \husson_get_h:n #1 {
  \tl_if_exist:cTF { g_husson_h_#1_tl }
    { \tl_use:c { g_husson_h_#1_tl } }
    { 0 }
}

\cs_new:Npn \husson_get_bar:n #1 {
  \tl_if_exist:cTF { g_husson_bar_#1_tl }
    { \tl_use:c { g_husson_bar_#1_tl } }
    { }
}

% Resolution de la marque active selon la premiere rubrique en vigueur sur la page
\cs_new:Npn \husson_render_bar: {
  \tl_set:Ne \l_tmpa_tl { \FirstMark{hussonnav} }
  \tl_if_empty:NF \l_tmpa_tl
    {
      \husson_get_bar:n { \l_tmpa_tl }
    }
}

% Composition du pied de page gauche (Auteur – Titre en italique)
\cs_new:Npn \husson_render_footer_left: {
  \tl_if_empty:NTF \g_husson_footerleft_tl
    {
      \tl_if_empty:NTF \g_husson_author_tl
        {
          \tl_if_empty:NF \g_husson_title_tl
            { \textit{ \tl_use:N \g_husson_title_tl } }
        }
        {
          \tl_use:N \g_husson_author_tl
          \tl_if_empty:NF \g_husson_title_tl
            { ~--~ \textit{ \tl_use:N \g_husson_title_tl } }
        }
    }
    {
      \tl_use:N \g_husson_footerleft_tl
    }
}

\NewDocumentCommand{\hussonnavstate}{mmm}{
  \husson_set_state:nnn {#1} {#2} {#3}
}

\NewDocumentCommand{\renderhussonnavbar}{}{
  \husson_render_bar:
}

\NewDocumentCommand{\hussongetfooterleft}{}{
  \husson_render_footer_left:
}

\NewDocumentCommand{\hussonnavmark}{m}{
  \InsertMark{hussonnav}{#1}
}
\ExplSyntaxOff

% Fallback automatique sur \@title et \@author si non definis par le filtre Lua
\makeatletter
\AddToHook{begindocument/before}{%
  \ExplSyntaxOn
  \tl_if_empty:NT \g_husson_title_tl
    {
      \tl_if_empty:NF \@title
        { \tl_gset:Nn \g_husson_title_tl { \@title } }
    }
  \tl_if_empty:NT \g_husson_author_tl
    {
      \tl_if_empty:NF \@author
        { \tl_gset:Nn \g_husson_author_tl { \@author } }
    }
  \ExplSyntaxOff
}
\makeatother

% Configuration du style de page avec scrlayer-scrpage
\clearpairofpagestyles
\KOMAoptions{headsepline=0.3pt, footsepline=0.3pt, footheight=18pt}
\setkomafont{headsepline}{\color{linegray}}
\setkomafont{footsepline}{\color{linegray}}
\ihead{\renderhussonnavbar}
\ifoot{\normalfont\scriptsize\textcolor{headingcharcoal}{\hussongetfooterleft}}
\cfoot[\pagemark]{}
\ofoot{\normalfont\scriptsize\textcolor{black}{\pagemark}}

% Desactivation propre sur la table des matieres
\AddToHook{cmd/tableofcontents/before}{%
  \pagestyle{plain}%
}
\AddToHook{cmd/tableofcontents/after}{%
  \pagestyle{scrheadings}%
}

\makeatletter
\@ifpackageloaded{caption}{}{\usepackage{caption}}
\AtBeginDocument{%
\ifdefined\contentsname
  \renewcommand*\contentsname{Table des matières}
\else
  \newcommand\contentsname{Table des matières}
\fi
\ifdefined\listfigurename
  \renewcommand*\listfigurename{Liste des Figures}
\else
  \newcommand\listfigurename{Liste des Figures}
\fi
\ifdefined\listtablename
  \renewcommand*\listtablename{Liste des tableaux}
\else
  \newcommand\listtablename{Liste des tableaux}
\fi
\ifdefined\figurename
  \renewcommand*\figurename{Figure}
\else
  \newcommand\figurename{Figure}
\fi
\ifdefined\tablename
  \renewcommand*\tablename{Tableau}
\else
  \newcommand\tablename{Tableau}
\fi
}
\@ifpackageloaded{float}{}{\usepackage{float}}
\floatstyle{ruled}
\@ifundefined{c@chapter}{\newfloat{codelisting}{h}{lop}}{\newfloat{codelisting}{h}{lop}[chapter]}
\floatname{codelisting}{Listing}
\newcommand*\listoflistings{\listof{codelisting}{Liste des Listings}}
\makeatother
\makeatletter
\makeatother
\makeatletter
\@ifpackageloaded{caption}{}{\usepackage{caption}}
\@ifpackageloaded{subcaption}{}{\usepackage{subcaption}}
\makeatother
\usepackage{bookmark}
\IfFileExists{xurl.sty}{\usepackage{xurl}}{} % add URL line breaks if available
\urlstyle{same}
\hypersetup{
  pdftitle={Validation de la barre de navigation},
  pdfauthor={Thomas Husson},
  pdflang={fr},
  colorlinks=true,
  linkcolor={NavyBlue},
  filecolor={Maroon},
  citecolor={Blue},
  urlcolor={NavyBlue},
  pdfcreator={LaTeX via pandoc}}


\title{Validation de la barre de navigation}
\usepackage{etoolbox}
\makeatletter
\providecommand{\subtitle}[1]{% add subtitle to \maketitle
  \apptocmd{\@title}{\par {\large #1 \par}}{}{}
}
\makeatother
\subtitle{Tests exhaustifs des niveaux H1 à H5, nav-title et transitions
de page}
\author{Thomas Husson}
\date{}
\begin{document}

\maketitle

\renewcommand*\contentsname{Table des matières}
{
\hypersetup{linkcolor=}
\setcounter{tocdepth}{3}
\tableofcontents
}


% Metadonnees du pied de page Husson
\hussondoctitle{Validation de la barre de navigation}
\hussondocauthor{Thomas Husson}

% Definitions des etats de la barre de navigation Husson
\hussonnavstate{1}{1}{\hussonnavfit{\hussonnavstandalone{Introduction}}}
\hussonnavstate{2}{2}{\hussonnavfit{\noindent\normalfont\scriptsize \textcolor{headingblue}{\textbf{MÉTHODES}}\hussonbarsep \textcolor{headingblue}{Design}}}
\hussonnavstate{3}{2}{\hussonnavfit{\noindent\normalfont\scriptsize \textcolor{headingblue}{\textbf{MÉTHODES}}\hussonbarsep \textcolor{headingblue}{Transporteurs}}}
\hussonnavstate{4}{2}{\hussonnavfit{\noindent\normalfont\scriptsize \textcolor{headingblue}{\textbf{MÉTHODES}}\hussonbarsep \textcolor{headingblue}{Perfusion}}}
\hussonnavstate{5}{2}{\hussonnavfit{\noindent\normalfont\scriptsize \textcolor{headingblue}{\textbf{MÉTHODES}}\hussonbarsep \textcolor{headingblue}{Biologie}\hussonsep \textcolor{headingblue}{Inflammation}\hussonsep \textcolor{headingblue}{Cytokines}\hussonsep \textbf{\textcolor{headingblue}{IL-6}}}}
\hussonnavstate{6}{2}{\hussonnavfit{\noindent\normalfont\scriptsize \textcolor{headingblue}{\textbf{MÉTHODES}}\hussonbarsep \textcolor{headingblue}{Biologie}\hussonsep \textbf{\textcolor{headingblue}{Hémolyse}}}}
\hussonnavstate{7}{2}{\hussonnavfit{\noindent\normalfont\scriptsize \textcolor{headingblue}{\textbf{MÉTHODES}}\hussonbarsep \textcolor{headingblue}{Statistiques}}}
\hussonnavstate{8}{8}{\hussonnavfit{\noindent\normalfont\scriptsize \textcolor{headingblue}{\textbf{RÉSULTATS}}\hussonbarsep \textcolor{headingblue}{Hémodynamique}}}
\hussonnavstate{9}{8}{\hussonnavfit{\noindent\normalfont\scriptsize \textcolor{headingblue}{\textbf{RÉSULTATS}}\hussonbarsep \textcolor{headingblue}{Gazométrie}}}
\hussonnavstate{10}{8}{\hussonnavfit{\noindent\normalfont\scriptsize \textcolor{headingblue}{\textbf{RÉSULTATS}}\hussonbarsep \textcolor{headingblue}{Métabolisme}}}
\hussonnavstate{11}{11}{\hussonnavfit{\hussonnavstandalone{Conclusion et perspectives}}}

\clearpage

\hussonnavmark{1}

\section*{Introduction}\label{introduction}
\addcontentsline{toc}{section}{Introduction}

Cette première section non numérotée sert de section liminaire ou
introductive. Elle permet de vérifier le comportement de la barre
lorsqu'un titre de niveau 1 ne possède aucune sous-section H2. La barre
doit afficher le titre seul à gauche avec son filet sobre.

\newpage

\clearpage

\section{Matériels et méthodes}\label{matuxe9riels-et-muxe9thodes}

\Needspace{8\baselineskip}

\hussonnavmark{2}

\subsection{Design expérimental et groupes
d'étude}\label{design-expuxe9rimental-et-groupes-duxe9tude}

Ce paragraphe décrit la méthodologie générale et la constitution des
cohortes. La barre de navigation doit afficher à ce stade :
\texttt{MÉTHODES\ \textbar{}\ Design\ \ \ Transporteurs\ \ \ Perfusion\ \ \ Biologie\ \ \ Statistiques}
avec \texttt{Design} en surbrillance active et souligné.

\vspace{5cm}

\begin{longtable}[]{@{}
  >{\raggedright\arraybackslash}p{(\linewidth - 6\tabcolsep) * \real{0.3400}}
  >{\centering\arraybackslash}p{(\linewidth - 6\tabcolsep) * \real{0.2404}}
  >{\centering\arraybackslash}p{(\linewidth - 6\tabcolsep) * \real{0.2334}}
  >{\centering\arraybackslash}p{(\linewidth - 6\tabcolsep) * \real{0.1862}}@{}}
\caption{Tableau comparatif des caractéristiques
initiales}\label{tbl-demo}\tabularnewline
\toprule\noalign{}
\begin{minipage}[b]{\linewidth}\raggedright
Paramètre
\end{minipage} & \begin{minipage}[b]{\linewidth}\centering
Groupe Contrôle
\end{minipage} & \begin{minipage}[b]{\linewidth}\centering
Groupe Traité
\end{minipage} & \begin{minipage}[b]{\linewidth}\centering
p-value
\end{minipage} \\
\midrule\noalign{}
\endfirsthead
\toprule\noalign{}
\begin{minipage}[b]{\linewidth}\raggedright
Paramètre
\end{minipage} & \begin{minipage}[b]{\linewidth}\centering
Groupe Contrôle
\end{minipage} & \begin{minipage}[b]{\linewidth}\centering
Groupe Traité
\end{minipage} & \begin{minipage}[b]{\linewidth}\centering
p-value
\end{minipage} \\
\midrule\noalign{}
\endhead
\bottomrule\noalign{}
\endlastfoot
Nombre de sujets & 25 & 25 & \textgreater{} 0.99 \\
Âge moyen (années) & 54.2 ± 6.1 & 55.0 ± 5.8 & 0.62 \\
Débit de perfusion (mL/min) & 120 ± 15 & 125 ± 12 & 0.18 \\
\end{longtable}

\newpage

\Needspace{8\baselineskip}

\hussonnavmark{3}

\subsection{Préparation des transporteurs
d'oxygène}\label{pruxe9paration-des-transporteurs-doxyguxe8ne}

Les transporteurs d'oxygène basés sur l'hémoglobine et les globules
rouges lavés ont été préparés selon les protocoles standardisés.

La barre de navigation doit maintenant afficher \texttt{Transporteurs}
en surbrillance active et souligné.

\newpage

\Needspace{8\baselineskip}

\hussonnavmark{4}

\subsection{Procédure de perfusion hépatique
normothermique}\label{procuxe9dure-de-perfusion-huxe9patique-normothermique}

Détails de la mise en route du circuit extracorporel et de la
canulation.

\newpage

\Needspace{8\baselineskip}

\subsection{Analyses biologiques des
perfusats}\label{analyses-biologiques-des-perfusats}

\Needspace{5\baselineskip}

\subsubsection{Évaluation de la réponse
inflammatoire}\label{uxe9valuation-de-la-ruxe9ponse-inflammatoire}

\Needspace{3\baselineskip}

\paragraph{Dosage des cytokines
circulantes}\label{dosage-des-cytokines-circulantes}

\Needspace{3\baselineskip}

\hussonnavmark{5}

\subparagraph{Interleukine 6}\label{interleukine-6}

Nous atteignons ici une branche profonde de niveau 5.

La position hiérarchique réelle est : - H1 : Matériels et méthodes
(\texttt{MÉTHODES}) - H2 : Analyses biologiques des perfusats
(\texttt{Biologie}) - H3 : Évaluation de la réponse inflammatoire
(\texttt{Inflammation}) - H4 : Dosage des cytokines circulantes
(\texttt{Cytokines}) - H5 : Interleukine 6 (\texttt{IL-6})

La barre doit afficher sur une seule ligne :
\texttt{MÉTHODES\ \textbar{}\ Design\ \ \ Transporteurs\ \ \ Perfusion\ \ \ Biologie\ ›\ Inflammation\ ›\ Cytokines\ ›\ IL-6\ \ \ Statistiques}

avec \texttt{Biologie\ ›\ Inflammation\ ›\ Cytokines\ ›\ IL-6} en bleu,
et \texttt{IL-6} en gras, le tout souligné par le trait bleu.

\newpage

\Needspace{5\baselineskip}

\hussonnavmark{6}

\subsubsection{Marqueurs d'hémolyse et souffrance
cellulaire}\label{marqueurs-dhuxe9molyse-et-souffrance-cellulaire}

Cette page teste le retour à un niveau supérieur (H3) sous le même H2
(\texttt{Biologie}). Les sous-branches précédentes
(\texttt{Inflammation}, \texttt{Cytokines}, \texttt{IL-6}) doivent être
complètement nettoyées.

La barre doit afficher :
\texttt{MÉTHODES\ \textbar{}\ Design\ \ \ Transporteurs\ \ \ Perfusion\ \ \ Biologie\ ›\ Hémolyse\ \ \ Statistiques}

\newpage

\Needspace{8\baselineskip}

\hussonnavmark{7}

\subsection{Analyses statistiques}\label{analyses-statistiques}

Les variables continues ont été exprimées en moyennes ± écart-type.

Cette section vérifie que le passage au dernier H2 du chapitre
réinitialise le déploiement de \texttt{Biologie} et active
\texttt{Statistiques}.

\newpage

\clearpage

\section{Résultats expérimentaux}\label{ruxe9sultats-expuxe9rimentaux}

\Needspace{8\baselineskip}

\hussonnavmark{8}

\subsection{Paramètres
hémodynamiques}\label{paramuxe8tres-huxe9modynamiques}

Premier H2 du second chapitre. La barre doit basculer sur le H1
\texttt{RÉSULTATS} et lister ses propres sous-sections H2.

\newpage

\Needspace{8\baselineskip}

\hussonnavmark{9}

\subsection{Gazométrie et
oxygénation}\label{gazomuxe9trie-et-oxyguxe9nation}

Texte de la première moitié de page sur la gazométrie.

\vspace{12cm}

\Needspace{8\baselineskip}

\hussonnavmark{10}

\subsection{Marqueurs métaboliques et
lactates}\label{marqueurs-muxe9taboliques-et-lactates}

Ce nouveau titre H2 commence au bas de la page courante. Selon la règle
validée du haut de page (\emph{Top-mark}), l'en-tête de cette page doit
continuer d'afficher \texttt{Gazométrie}. Le titre \texttt{Métabolisme}
prendra effet dans l'en-tête de la page suivante.

\newpage

Suite de la section sur les marqueurs métaboliques et la clairance des
lactates sur la page suivante. L'en-tête doit ici afficher
\texttt{Métabolisme} actif.

\newpage

\clearpage

\hussonnavmark{11}

\section*{Conclusion et perspectives}\label{conclusion-et-perspectives}
\addcontentsline{toc}{section}{Conclusion et perspectives}

Cette section finale n'a pas de sous-section H2. Elle doit s'afficher
sobrement comme \texttt{Conclusion\ et\ perspectives} à gauche, avec la
ligne fine en-dessous.




\end{document}
